% =============================================================================
% ECSS-E-ST-10-06C SECTION <6>: ENVIRONMENT AND CONSTRAINTS DESCRIPTION
% =============================================================================
\subsection{6. Environment and Constraints Description}

This section characterizes the natural space environment, induced launch and contamination environments, internal micro-vibration disturbances, and programmatic, security, and regulatory constraints that bound Configuration Item \textbf{\texttt{CI-03}: Space Segment (\abbr{SSEG}, \abbr{WBS} Element \texttt{123456.03})} across all life-profile phases defined in Section~5. Detailed engineering design trade-offs and shielding justifications are documented in the Space Segment Design Justification File (\texttt{EWAS-HSB-SSEG-DDJF-03}).

\subsubsection{6.1 Space Radiation and Plasma Environment (1,000~km LEO vs.\ 35,786~km GEO, \texttt{ECSS-E-ST-10-04C})}

Because the \abbr{EWAS} Space Segment operates simultaneously in circular Low Earth Orbit (1,000~km, $80^\circ$ inclination for \texttt{SATL}) and equatorial Geostationary Earth Orbit (35,786~km, $0^\circ$ inclination for \texttt{SATG}), its two flight layers encounter distinct ionizing radiation, non-ionizing particle, and plasma charging environments (\textbf{ECSS-E-ST-10-04C} Clause 9, evaluated using AP-8/AE-8, IRENE AE9/AP9, and SAPPHIRE solar particle models at a $95\%$ confidence level):
\begin{enumerate}
    \item \textbf{1,000~km, $80^\circ$ Inclination \abbr{LEO} Regime (\texttt{SATL}, 7-Year Design Lifetime \texttt{R-016}):}
    \begin{itemize}
        \item \emph{Trapped Protons and South Atlantic Anomaly (\abbr{SAA}) Transits:} At 1,000~km altitude, the constellation orbits below the main proton peak of the inner Van Allen belt (which rises steeply above 1,150--1,200~km). However, over the South Atlantic Anomaly (\abbr{SAA}, $15^\circ\text{S}\text{--}45^\circ\text{S}$, $90^\circ\text{W}\text{--}0^\circ\text{E}$), each \texttt{SATL} node traverses intense trapped proton fluxes ($E_p > 10\text{--}100\text{ MeV}$, peak omnidirectional flux $\sim 10^3\text{--}10^4\text{ protons}/(\text{cm}^2\cdot\text{s})$) during 4--6 orbits per day ($\sim 8\text{--}12\text{ min}$ per transit). In addition, the $80^\circ$ orbital inclination exposes the spacecraft during polar cap crossings ($> 60^\circ\text{N/S}$, where geomagnetic cutoff rigidity approaches zero) to unshielded Solar Energetic Particles (\abbr{SEP}) and Galactic Cosmic Rays (\abbr{GCR}).
        \item \emph{Total Ionizing Dose (\abbr{TID}) and Displacement Damage (\abbr{DD}) Envelope:} Over the 7-year mission lifetime (\texttt{R-016}), including a mandatory Radiation Design Margin (\abbr{RDM}) of $2\times$ per \textbf{ECSS-E-ST-10-04C} and \textbf{ECSS-E-ST-20C} (\textbf{AD-07}), the cumulative dose behind a reference 4~mm aluminium boundary is $\sim 18\text{--}30\text{ krad}(\mathrm{Si})$, constraining internal shielded electronics locations to $\le 10\text{--}15\text{ krad}(\mathrm{Si})$. Cumulative Non-Ionizing Energy Loss (\abbr{NIEL} / Displacement Damage equivalent fluence of $\sim 2 \times 10^{10}\text{ protons}/\text{cm}^2$ at $10\text{ MeV}$) imposes dark-current stability constraints on the $\mathrm{HgCdTe}$ \abbr{SWIR}/\abbr{MWIR} Focal Plane Arrays (\abbr{FPA}s) and star-tracker CMOS detectors.
        \item \emph{Single-Event Effects (\abbr{SEE}) and Focal-Plane Proton Flux:} High-energy protons in the \abbr{SAA} and heavy cosmic-ray ions at high latitudes induce Single-Event Upsets (\abbr{SEU}), Single-Event Functional Interrupts (\abbr{SEFI}), Single-Event Latch-ups (\abbr{SEL}, requiring immunity or protection up to $\text{LET}_{\text{th}} \ge 37\text{ MeV}\cdot\text{cm}^2/\text{mg}$), and transient proton ionization spikes on the infrared detector pixels (\texttt{R-019}, \texttt{R-022}).
    \end{itemize}
    \item \textbf{35,786~km Equatorial \abbr{GEO} Regime (\texttt{SATG}, 15-Year Design Lifetime \texttt{R-016}):}
    \begin{itemize}
        \item \emph{Outer-Belt Relativistic Electrons and Solar Proton Events:} At geostationary altitude, trapped proton flux is negligible above $1\text{ MeV}$, and the radiation environment is dominated by outer Van Allen belt relativistic electrons ($E_e > 0.5\text{--}5.0\text{ MeV}$), bremsstrahlung X-rays, unattenuated Solar Proton Events (\abbr{SPE}), and cosmic-ray heavy ions. Over the 15-year operational lifetime (\texttt{R-016}, including a $2\times$ \abbr{RDM}), the cumulative \abbr{TID} behind a 4~mm aluminium reference boundary reaches $\sim 60\text{--}100\text{ krad}(\mathrm{Si})$, constraining internal vault-shielded locations to $\le 20\text{--}25\text{ krad}(\mathrm{Si})$.
        \item \emph{Surface and Internal Dielectric Charging Environment:} During geomagnetic substorms in the Earth's magnetotail, \texttt{SATG} is immersed in hot, tenuous magnetospheric plasma ($k_B T_e \approx 1\text{--}20\text{ keV}$) while energetic electrons ($> 100\text{ keV}$) penetrate into dielectrics, imposing strict surface and internal electrostatic discharge (\abbr{ESD}) limits ($< 100\text{ V}$ differential potential) per \textbf{ECSS-E-ST-20-06C} (\textbf{AD-07}).
    \end{itemize}
\end{enumerate}

\subsubsection{6.2 Thermal, Solar Illumination, and Eclipse Environment (\texttt{ECSS-E-ST-31C})}

The space thermal environment (\textbf{ECSS-E-ST-31C}, \textbf{AD-08}; \textbf{ECSS-E-ST-10-04C} Clause 6, \textbf{AD-06}) subjects the two constellation tiers to distinct radiative heat flux extremes and orbital eclipse regimes:
\begin{itemize}
    \item \textbf{External Radiative Heat Fluxes:} Both constellations are exposed to direct solar irradiance varying seasonally from $S_{\min} = 1{,}322\text{ W/m}^2$ (aphelion) to $S_{\max} = 1{,}414\text{ W/m}^2$ (perihelion), with a nominal solar constant of $S_0 = 1{,}361\text{ W/m}^2$. At 1,000~km \abbr{LEO}, \texttt{SATL} additionally receives Earth-reflected solar albedo ($a = 0.25\text{--}0.38$, corresponding to $150\text{--}380\text{ W/m}^2$ depending on cloud and ice cover over the Baltic and Arctic) and Earth outgoing longwave infrared radiation ($q_{\text{EIR}} = 210\text{--}260\text{ W/m}^2$). At 35,786~km \abbr{GEO}, Earth albedo and planetary infrared fluxes on \texttt{SATG} are $< 7\text{ W/m}^2$.
    \item \textbf{Eclipse Regimes and Thermal Cycling Constraints:}
    \begin{enumerate}
        \item \emph{\texttt{SATL} (1,000~km, $80^\circ$ \abbr{LEO}):} Completes $\sim 13.7$ orbits per day ($a = 7{,}378.14\text{ km}$, period $T = 105.12\text{ min}$). When the solar beta angle approaches $\beta_{\odot} \approx 0^\circ$, the maximum cylindrical Earth-shadow eclipse duration per orbit is given by:
        \[
            \Delta t_{\text{ecl,max}} = \frac{T}{\pi}\arcsin\!\left(\frac{R_E}{a}\right) = \frac{105.12\text{ min}}{180^\circ} \times 59.82^\circ = 34.93\text{ min} \approx 35.0\text{ min}
        \]
        Over its 7-year lifetime, each \texttt{SATL} node experiences \textbf{25,000 to 35,000 sunlit-to-eclipse thermal cycles}, imposing severe thermo-mechanical fatigue constraints on solar array interconnects, $\mathrm{Li\text{-}ion}$ battery depth of discharge ($\text{DoD} \le 30\%$), and star-tracker-to-telescope thermo-elastic alignment stability ($\le 10\text{--}15~\mu\mathrm{rad}$ \abbr{AKE}).
        \item \emph{\texttt{SATG} (35,786~km \abbr{GEO}):} Remains in continuous sunlight throughout most of the year, interrupted by two 45-day eclipse seasons centered on the vernal and autumnal equinoxes ($\sim 90$ eclipses/year; $\sim 1{,}350$ eclipses over 15~years) with a maximum daily eclipse duration of \textbf{71.5~minutes} ($\text{DoD} \le 65\%$).
    \end{enumerate}
    \item \textbf{Dual-Regime Temperature Boundary Constraints:} Each spacecraft must simultaneously sustain two widely separated temperature zones: (1)~the warm electronics bus ($263\text{--}323\text{ K}$) and $\mathrm{Li\text{-}ion}$ batteries ($273\text{--}303\text{ K}$), and (2)~the cryogenic telescope optics ($180\text{--}210\text{ K}$) and $\mathrm{HgCdTe}$ \abbr{SWIR}/\abbr{MWIR} focal plane arrays held at \textbf{$65\text{ K}\text{--}95\text{ K}$} within $\pm 0.1\text{ K}$.
\end{itemize}

\subsubsection{6.3 Vacuum, Residual Atmosphere, Outgassing, and Contamination Constraints}

\begin{itemize}
    \item \textbf{Hard Vacuum and Cryogenic Molecular Contamination (\texttt{ECSS-Q-ST-70-02C}):} Ambient pressure is $< 10^{-9}\text{ hPa}$ at 1,000~km \abbr{LEO} and $< 10^{-11}\text{ hPa}$ at 35,786~km \abbr{GEO}. In hard vacuum, cryogenic focal planes ($65\text{--}95\text{ K}$) and cold optical mirrors ($180\text{--}210\text{ K}$) act as molecular cold traps where condensing water ice ($\mathrm{H_2O}$) or hydrocarbon films of $10\text{--}50\text{ nm}$ thickness absorb strongly in the $2.7~\mu\mathrm{m}$ (\abbr{SWIR}) and $4.2\text{--}4.4~\mu\mathrm{m}$ (\abbr{MWIR}) bands and degrade stray-light rejection (\texttt{R-009}, \texttt{R-025}). All non-metallic materials are therefore constrained to Total Mass Loss $\mathrm{TML} < 1.0\%$ and Collected Volatile Condensable Material $\mathrm{CVCM} < 0.01\%$, and propulsion plumes must remain clear of optical apertures.
    \item \textbf{Exospheric Drag and Atomic Oxygen (\abbr{AO}) at 1,000~km:} At 1,000~km altitude, the residual exosphere ($\rho \approx 10^{-15}\text{ to }10^{-14}\text{ kg/m}^3$ across the 11-year solar cycle) produces low aerodynamic drag (semi-major axis decay of $\sim 1\text{--}5\text{ m/year}$) and low Atomic Oxygen fluence ($< 10^{18}\text{ atoms/cm}^2\text{/year}$). While differential drag requires periodic station-keeping burns (Mode N4), the negligible natural decay rate strictly mandates active retrograde deorbit propulsion at End of Life (\texttt{R-018}).
\end{itemize}

\subsubsection{6.4 Micrometeoroids and Orbital Debris (MMOD) and Gravitational Perturbations}

\begin{itemize}
    \item \textbf{MMOD Spatial Density and Collision Risk Constraint (\texttt{ECSS-E-ST-10-04C} Clause 10):} The $800\text{--}1{,}050\text{ km}$ high-inclination \abbr{LEO} band exhibits a high spatial density of catalogued and uncatalogued orbital debris (modeled via ESA \texttt{MASTER-8} at hypervelocity impact speeds averaging $v_{\text{rel}} \approx 11\text{--}14.5\text{ km/s}$). This environment imposes a binding spacecraft-level Probability of No Penetration constraint of $\text{PNP} \ge 0.99$ over the design lifetime against untrackable particles ($< 1.0\text{ cm}$) and mandates active Collision Avoidance Manoeuvres (\abbr{CAM}) whenever Ground Segment Conjunction Data Messages (\texttt{IF-EXT-04}) indicate a collision probability $P_c > 10^{-4}$.
    \item \textbf{Gravitational, Geopotential, and Solar Radiation Pressure Perturbations:} At 1,000~km, $80^\circ$ inclination, Earth's oblateness ($J_2 = 1.08263 \times 10^{-3}$) induces a secular nodal precession of $\dot{\Omega} \approx -1.04^\circ/\text{day}$ and apsidal rotation of $\dot{\omega} \approx +2.85^\circ/\text{day}$ ($\Delta v_{\text{SK}} \le 4\text{--}6\text{ m/s/year}$ for differential phase maintenance). At 35,786~km \abbr{GEO}, lunisolar gravitational attraction ($\sim 45\text{--}51\text{ m/s/year}$ N--S inclination drift), Earth triaxiality ($J_{22}$ E--W drift $\sim 2\text{ m/s/year}$), and solar radiation pressure govern the \texttt{SATG} station-keeping constraint.
\end{itemize}

\subsubsection{6.5 Mechanical, Vibro-Acoustic, and In-Orbit Micro-Vibration Constraints (\texttt{ECSS-E-ST-32C})}

\begin{itemize}
    \item \textbf{Launch Quasi-Static, Vibration, Acoustic, and Shock Envelope (\textbf{AD-09}):} During launch ascent (\texttt{IF-EXT-02}), both \texttt{SATL} (mounted 9-per-stack on a cylindrical multi-satellite dispenser) and \texttt{SATG} must withstand quasi-static axial accelerations up to $6.0\,g$ and lateral accelerations up to $3.0\,g$, sinusoidal vibrations ($5\text{--}100\text{ Hz}$, up to $1.5\,g$), random vibro-acoustic loads ($20\text{--}2{,}000\text{ Hz}$, $7\text{--}11\,g_{\text{rms}}$; acoustic $138\text{--}142\text{ dB}$ OASPL), and separation shock ($1{,}500\text{--}2{,}000\,g$ SRS above $1{,}000\text{ Hz}$). Stowed fundamental structural frequencies are constrained to $f_{\text{axial}} \ge 35\text{ Hz}$ and $f_{\text{lateral}} \ge 20\text{ Hz}$ for \texttt{SATL} ($f_{\text{axial}} \ge 31\text{ Hz}$, $f_{\text{lateral}} \ge 15\text{ Hz}$ for \texttt{SATG}).
    \item \textbf{In-Orbit Micro-Vibration (Jitter) Constraint (\texttt{R-024}):} During nominal operations, internal mechanical disturbances are generated by reaction-wheel imbalances ($10\text{--}300\text{ Hz}$), pulse-tube cryocooler compressor strokes ($45\text{--}85\text{ Hz}$), and dynamic slewing of the \abbr{SWIR}/\abbr{MWIR} sensor gimbals/mirrors ($\ge 15^\circ/\text{s}$), directional \abbr{RF} antennas, and agile laser \abbr{OISL} terminals. These combined disturbances are constrained by the requirement to maintain optical Line-of-Sight (\abbr{LOS}) pointing stability strictly below $< 2.0\text{ arcsec/s}$ ($\le 9.7~\mu\mathrm{rad/s}$ RMS, \texttt{R-024}) and laser \abbr{OISL} beam pointing below $< 1.5~\mu\mathrm{rad}$.
\end{itemize}

\subsubsection{6.6 Programmatic, Operational, Security, and Regulatory Constraints}

In addition to physical environments, the \abbr{EWAS} Space Segment is bounded by three system-level constraint categories flowed down from \textbf{AD-01}:
\begin{enumerate}
    \item \textbf{Non-Offensive Civil-Protection and Functional Scope Constraints (\texttt{R-001}, \texttt{R-002}, \texttt{R-048}, \texttt{R-049}):} The Space Segment is strictly constrained to defensive civil-protection early warning---detecting missile launches globally and regionally (\texttt{R-049}), tracking trajectories stereoscopically, and downlinking observation packages and preliminary trajectory estimates to the Ground Segment (\texttt{R-023}, \texttt{R-045}). It is prohibited from incorporating weapon-targeting, fire-control, or kinetic-interceptor guidance functions (\texttt{R-001}), and is not required to perform warhead yield characterization or post-impact damage assessment (\texttt{R-002}).
    \item \textbf{Sovereign Data Security, Archival Governance, and Electronic Warfare Resilience (\texttt{R-003}, \texttt{R-004}, \texttt{R-031}, \texttt{R-032}, \texttt{R-035}--\texttt{R-037}):} All space-to-ground and inter-satellite links must enforce hardware bulk encryption and cryptographic telecommand authentication (\texttt{R-031}), accompanied by spread-spectrum \abbr{RF} and narrow-beam $1{,}550\text{ nm}$ optical anti-jamming/anti-spoofing protection (\texttt{R-032}). Downlinked telemetry and sensor observation records are formatted for 5-year sovereign archival in the Ground Segment (\texttt{R-036}) and restricted exclusively to calibration, training, verification, and algorithm refinement (\texttt{R-004}).
    \item \textbf{International Regulatory and Space Sustainability Constraints (\texttt{R-018}, \textbf{AD-12}, \textbf{RD-01}, \textbf{RD-03}):} All S-band, X-band, and Ka-band \abbr{RF} links must comply with International Telecommunication Union (\abbr{ITU}) Radio Regulations (\textbf{RD-03}) spectral masks and power-flux-density limits, while laser links incorporate aviation and space safety exclusion interlocks. Every \texttt{SATL} and \texttt{SATG} spacecraft is bound by mandatory active End-of-Life disposal and passivation (\texttt{R-018}, \textbf{ECSS-U-AS-10C} \textbf{AD-12}, and \textbf{ESA Zero Debris Policy} \textbf{RD-01}).
\end{enumerate}

Table~\ref{tab:sseg_env_constraints} summarizes the quantitative environmental regimes and governing boundary constraints for the \abbr{EWAS} Space Segment.

\begin{table}[H]
\centering
\caption{Summary of EWAS Space Segment Environmental Regimes and Governing Constraints}
\label{tab:sseg_env_constraints}
{\scriptsize\setlength{\tabcolsep}{4pt}\renewcommand{\arraystretch}{1.05}%
\arrayrulecolor{TableBorderBlue}%
\rowcolors{2}{white}{TableRowLightBlue}%
\noindent
\begin{tabularx}{\textwidth}{|>{\raggedright\arraybackslash}p{2.4cm}|>{\raggedright\arraybackslash}p{2.8cm}|>{\raggedright\arraybackslash}p{2.8cm}|>{\raggedright\arraybackslash}X|}
    \hline
    \rowcolor{TableHeaderBlue}
    \thead{Environment / Constraint} & \thead{\texttt{SATL} Envelope (1,000~km, $80^\circ$ LEO)} & \thead{\texttt{SATG} Envelope (35,786~km, $0^\circ$ GEO)} & \thead{Governing Environmental Constraint \& Applicable Standard} \\ \hline
    \textbf{1. Ionizing Radiation (\abbr{TID}, $2\times$ RDM)} & 18--30~krad($\mathrm{Si}$) (4~mm Al); $\le 10\text{--}15$~krad($\mathrm{Si}$) internal (7~yr, \texttt{R-016}) & 60--100~krad($\mathrm{Si}$) (4~mm Al); $\le 20\text{--}25$~krad($\mathrm{Si}$) internal (15~yr, \texttt{R-016}) & Cumulative ionizing dose and displacement damage survival over $\ge 7$~yr (\abbr{LEO}) and $\ge 15$~yr (\abbr{GEO}) with $2\times$ \abbr{RDM} (\texttt{R-016}, \textbf{ECSS-E-ST-10-04C} \textbf{AD-06}, \textbf{AD-07}). \\ \hline
    \textbf{2. Protons, \abbr{SEE} \& Plasma Charging} & \abbr{SAA} protons ($10^3\text{--}10^4\text{ p/cm}^2\text{/s}$); polar \abbr{SEP}/\abbr{GCR}; \abbr{FPA} proton spikes & Outer-belt relativistic $e^-$ ($> 0.5\text{--}5\text{ MeV}$); substorm charging ($1\text{--}20\text{ keV}$) & Destructive latch-up immunity ($\text{LET}_{\text{th}} \ge 37\text{ MeV}\cdot\text{cm}^2\text{/mg}$); single-frame \abbr{FPA} proton despiking (\texttt{R-019}, \texttt{R-022}); surface potential $< 100\text{ V}$ (\textbf{ECSS-E-ST-20-06C}). \\ \hline
    \textbf{3. Thermal, Solar \& Eclipse Cycling} & $S_0 = 1{,}361\text{ W/m}^2$; albedo $0.25\text{--}0.38$; $\le 35.0\text{ min}$ eclipse/orbit ($\sim 30{,}000$ cycles/7~yr) & $S_0 = 1{,}361\text{ W/m}^2$; two 45-day equinox seasons/yr ($\le 71.5\text{ min}$ max eclipse) & Cryogenic \abbr{FPA} stability at $65\text{--}95\text{ K} \pm 0.1\text{ K}$; bus electronics $-10^\circ\text{C}$ to $+50^\circ\text{C}$; thermo-elastic $\text{AKE} \le 10\text{--}15~\mu\mathrm{rad}$ (\textbf{ECSS-E-ST-31C} \textbf{AD-08}). \\ \hline
    \textbf{4. Vacuum, Outgassing \& Drag} & $< 10^{-9}\text{ hPa}$; low drag ($\rho \sim 10^{-15}\text{--}10^{-14}\text{ kg/m}^3$) & $< 10^{-11}\text{ hPa}$; hard interplanetary-class vacuum & Material outgassing limits $\mathrm{TML} < 1.0\%$, $\mathrm{CVCM} < 0.01\%$ (\textbf{ECSS-Q-ST-70-02C}); zero plume impingement on \texttt{R-025} baffles and \abbr{OISL} optics. \\ \hline
    \textbf{5. Micrometeoroids \& Debris (\abbr{MMOD})} & High spatial density in 800--1,050~km polar band ($v_{\text{rel}} \approx 11\text{--}14.5\text{ km/s}$) & Low spatial density; low relative velocity ($\sim 0.5\text{--}1.5\text{ km/s}$) & Spacecraft Probability of No Penetration $\text{PNP} \ge 0.99$ over design life; mandatory Collision Avoidance Manoeuvre for $P_c > 10^{-4}$ (\texttt{IF-EXT-04}, \textbf{AD-06}). \\ \hline
    \textbf{6. Launch Dynamic \& Vibro-Acoustic Loads} & Quasi-static $\le 6.0\,g$ axial, $\le 3.0\,g$ lat.; $7\text{--}11\,g_{\text{rms}}$ random; 9-sat dispenser stack & Quasi-static $\le 6.0\,g$ axial, $\le 2.5\,g$ lat.; acoustic $138\text{--}142\text{ dB}$; shock $\le 2{,}000\,g$ & Stowed structural stiffness $f_{\text{ax}} \ge 35\text{ Hz}$, $f_{\text{lat}} \ge 20\text{ Hz}$ (\texttt{SATL}) and $f_{\text{ax}} \ge 31\text{ Hz}$, $f_{\text{lat}} \ge 15\text{ Hz}$ (\texttt{SATG}) per \textbf{ECSS-E-ST-10-03C}/\textbf{32C} (\textbf{AD-05}, \textbf{AD-09}). \\ \hline
    \textbf{7. In-Orbit Micro-Vibration (Jitter)} & Wheels ($10\text{--}300\text{ Hz}$), cooler ($45\text{--}85\text{ Hz}$), slewing sensors, antennas \& \abbr{OISL} & Wheels, dual coolers, \& cross-layer antenna/\abbr{OISL} tracking mechanisms & Coupled optical \abbr{LOS} jitter/drift strictly $< 2.0\text{ arcsec/s}$ ($\le 9.7~\mu\mathrm{rad/s}$ RMS, \texttt{R-024}) and $< 1.5~\mu\mathrm{rad}$ \abbr{OISL} beam pointing stability. \\ \hline
    \textbf{8. Security, Scope \& Sustainability} & Non-offensive civil protection (\texttt{R-001}); contested EW environment; $< 5\text{ yr}$ \abbr{EOL} demise & Non-offensive (\texttt{R-001}); contested EW environment; $\Delta H \ge +300\text{ km}$ graveyard & Hardware \abbr{COMSEC} (\texttt{R-031}); anti-jamming (\texttt{R-032}); exclusion of weapon guidance (\texttt{R-001}); active \abbr{EOL} disposal \& passivation (\texttt{R-018}, \textbf{AD-12}). \\ \hline
\end{tabularx}}
\end{table}
